% Canonical Paper I source.

This chapter gives the ``hole'' and the ``cap'' a finite geometric carrier.  A
decorated directed acyclic network is thickened to a three-dimensional body
whose side surface carries the holes and whose interior carries the interface
disks; cutting the body along a disk exposes two caps.

\subsection{The decorated ball-and-tube thickening}

\begin{definition}[Decorated ball-and-tube thickening]
\label{def:decorated-ball-tube-thickening}
Let $G$ be a finite directed acyclic network with ordered typed external inputs
and outputs, ordered input/output slots and deterministic partial functions at
each gate, immutable sources, distinct occurrences, explicit copy/discard, and
binding records.  Forget the direction to obtain the underlying multigraph
$|G|$.  Choose, for each gate $v$, an oriented $3$-ball $B_v$ with disjoint
parametrized port disks on its boundary, one per slot; for each internal wire
$e$ choose a tube $T_e=D^2\times[0,1]$ and glue its end disks to the matching
port disks, orientation-compatibly.  The resulting body is $H_G$; the external
port disks are $E_G$, and the observed side surface is
\[
  S_G=\partial H_G\setminus\bigcup_{D\in E_G}\operatorname{int}(D).
\]
\end{definition}

This is a declared minimal thickening rule, not an intrinsic geometry forced by
the operation arity; the bare $H_G$ or $S_G$ carries no program identity.

\begin{proposition}[Local pants and no-hidden-handle genus]
\label{prop:local-pants-handlebody-genus}
\begin{enumerate}
  \item A gate with total port count three has its sphere minus the three open
  port disks a \emph{pair of pants}; total port count two gives a cylinder, one
  gives a disk.  The $2\to1$ and $1\to2$ decorations are distinct.
  \item For a connected $G$, $H_G$ is a handlebody with
  \[
    g(H_G)=\beta_1(|G|),
    \qquad
    \chi(S_G)=2-2\beta_1(|G|)-b,
  \]
  where $b=|E_G|$ is the number of external ports.
\end{enumerate}
\end{proposition}

\begin{proof}
(1) is the classification of surfaces with boundary.  (2) Assemble along a
spanning tree of $|G|$: each gate ball is a $0$-handle, each tree edge glues two
balls into one, and each remaining internal edge adds a $1$-handle; hence the
genus is the number of independent cycles.  The closed boundary of a genus-$g$
handlebody has Euler characteristic $2-2g$, and deleting $b$ open port disks
lowers it by $b$.
\end{proof}

\subsection{Interface disks and the hole--disk complementarity}

\begin{proposition}[Pants gluing along an interface disk]
\label{prop:pants-gluing-interface-disk}
Let two three-port gate balls be connected by one output--input tube.  At the
mid-section take the disk $D_e=D^2\times\{1/2\}$ and its boundary circle
$C_e=\partial D_e$.  Cutting along $D_e$ and absorbing the half-tubes into the
two collars, the two bodies are $3$-balls and the side surfaces are pants;
regluing along $D_e$ is, on the side surfaces, gluing along $C_e$:
\[
  B_1\cup_{D_e}B_2\cong B^3,
  \qquad
  S_G\cong P_1\cup_{C_e}P_2,
  \qquad
  \partial D_e=C_e .
\]
\end{proposition}

\begin{proof}
Two balls joined by one $1$-handle is again a ball (the underlying graph is a
tree); the interior disk is not part of the exposed boundary, and its circle
becomes the seam of the remaining boundary, a four-holed sphere with
$\chi=-1-1=-2$.
\end{proof}

\begin{theorem}[Hole--disk complementarity]
\label{thm:hole-disk-complementarity}
In the four-holed side surface $S_G$ and the thickened body $H_G$ of the
preceding proposition:
\begin{enumerate}
  \item \textbf{Surface non-fillability (hole).} $[C_e]\neq0\in H_1(S_G;\mathbb F_2)$:
  $C_e$ bounds no disk inside $S_G$.
  \item \textbf{Body fillability (disk/cap).} Under the inclusion
  $i:S_G\hookrightarrow H_G$, $i_*[C_e]=0\in H_1(H_G;\mathbb F_2)$, witnessed by
  the explicit disk $D_e\subset H_G$ with $\partial D_e=C_e$.
  \item \textbf{Cutting exposes caps.} Cutting $H_G$ along $D_e$ splits the single
  internal disk into two copies $D_e^+,D_e^-$ with opposite boundary
  orientations, each capping one pants opening.
\end{enumerate}
Hence the surface gives the hole its meaning and the body gives the disk its
meaning, the two being linked by the oriented common edge $C_e=\partial D_e$.
\end{theorem}

\begin{proof}
(1) and (2) are the mod-$2$ homology computation recorded and verified in
Appendix~\ref{app:pants-cellulation-verification}: the side-surface boundary
matrix has rank $110$, rising to $111$ on adjoining $C_e$, so $C_e$ is not a
boundary in $S_G$, while $C_e=\partial D_e$ is a boundary in $H_G$.  (3) is the
cut description of Proposition~\ref{prop:pants-gluing-interface-disk}.
\end{proof}

\begin{remark}[Down-set frontier disks]
\label{rem:downset-frontier-disks}
For a down-set $U\in\mathcal O(P)$ of the dependency poset of
Theorem~\ref{thm:birkhoff-downset-spinal}, cutting every internal tube that
points from $U$ to $P\setminus U$ yields the family of interface disks
$\{D_e\}_{e\in\partial^+U}$ representing the computational frontier.
\end{remark}

\subsection{The seam-circle holonomy}

The hole--disk complementarity of
Theorem~\ref{thm:hole-disk-complementarity} is topological.  It has an analytic
counterpart: the seam circle carries a holonomy of the surface potential, which
is the boundary reading of the bulk curvature of
Section~\ref{sec:differential-theory}.

\begin{definition}[Surface potential]
\label{def:surface-potential}
Let the arithmetic data $(a,v)$ be declared along the side surface $S_G$, with
$a$ the assignment and $v$ the multiplicative coordinate.  The \emph{surface
potential} is the one-form
\[
  \beta:=(-\lambda a\,dv)\big|_{S_G},
\]
the restriction to $S_G$ of the primitive $-\lambda a\,dv$ whose exterior
derivative is the contact curvature $d\alpha=-\lambda\,da\wedge dv$.
\end{definition}

\begin{proposition}[Seam-circle holonomy]
\label{prop:seam-circle-holonomy}
For the seam circle $C_e=\partial D_e$ of
Proposition~\ref{prop:pants-gluing-interface-disk}, assume that the declared
data $a,v$ extend smoothly to a neighborhood of $D_e$ and that $\lambda$ is
constant.  Write $\widetilde\beta=-\lambda a\,dv$ on that neighborhood, with
boundary restriction $\beta$.  With the induced boundary orientation, the loop integral
\[
  \oint_{C_e}\beta=\oint_{C_e}(-\lambda a\,dv)
\]
is the \emph{seam-circle holonomy}.  By the holographic Stokes theorem of
Proposition~\ref{prop:holographic-stokes-minimal}, this holonomy is the boundary
reading of the bulk curvature $\int_\Sigma d\alpha$ over the disk $\Sigma=D_e$
that fills $C_e$; hence the same circle that is topologically non-fillable in
$S_G$ ($[C_e]\neq0$) is analytically the carrier of a holonomy that the filling
disk realizes.
\end{proposition}

\begin{proof}
The first assertion is the definition.  The second is
Proposition~\ref{prop:holographic-stokes-minimal} applied to the filling disk
$\Sigma=D_e$ with boundary $\partial\Sigma=C_e$: the exactness of $d\alpha$ makes
the bulk area integral equal to the boundary holonomy.  The juxtaposition with
the topological non-fillability is Theorem~\ref{thm:hole-disk-complementarity}.
\end{proof}

\begin{remark}[The analytic thickening is a declared structure]
\label{rem:analytic-thickening-declared}
The surface potential $\beta$ is well-defined only after the arithmetic data
$(a,v)$ are declared on $S_G$; the ball-and-tube thickening of
Definition~\ref{def:decorated-ball-tube-thickening} is purely topological and
does not by itself carry $a,v$.  Extending the AES data from the bulk to the
thickened body and restricting to $S_G$ is the \emph{analytic thickening}, a
declared structure not yet constructed in general; the holonomy of
Proposition~\ref{prop:seam-circle-holonomy} is meaningful precisely under that
declaration.  This is the analytic counterpart of the observation that a bare
$S_G$ carries no program identity.
\end{remark}

\subsection{Drilling a tunnel and adding a surface handle}
\label{subsec:tunnel-handle-surgery}

The preceding construction distinguishes a circle that is unfillable on the
side surface from a disk that fills it in the body.  Drilling through the body
can change that filling relation.  We first describe the topological operation,
keeping all external port disks fixed; no arithmetic assignment is assumed.

\begin{proposition}[Two mouths and one new handle]
\label{prop:tunnel-boundary-surgery}
Let $H$ be a compact oriented $3$-manifold, let $\Sigma$ be a connected component
of $\partial H$, and mark $b$ disjoint external port disks on $\Sigma$.  Write
$S=\Sigma\setminus\bigcup_j\operatorname{int}E_j\cong S_{g,b}$.
Let a properly embedded arc have a closed tubular neighborhood
$N\cong D^2\times I$, meeting $\partial H$ exactly in its two end disks
$D_0,D_1\subset\operatorname{int}S$.  Drill out the relative interior of $N$
and round the corners.  The replacement side surface is
\[
 S'=\bigl(S\setminus(\operatorname{int}D_0\cup\operatorname{int}D_1)\bigr)
       \cup_{\partial D_0\sqcup\partial D_1}(S^1\times I),
 \qquad S'\cong S_{g+1,b}.
\]
Thus the tunnel wall becomes part of the boundary, the number of external
ports stays $b$, and $\chi(S')=\chi(S)-2$.  Attaching an oriented $3$-dimensional
$1$-handle along $D_0,D_1$ induces the same abstract surface replacement.
\end{proposition}

\begin{proof}
The two end disks cease to be boundary faces of the remaining body, whereas
the lateral boundary of $N$ becomes exposed.  Removing the two open disks
preserves connectedness and decreases Euler characteristic by two.  The
inserted annulus and its attaching circles all have Euler characteristic zero,
so they do not change that number.  The result is orientable, connected, and
has the original $b$ boundary circles.  The formula $\chi=2-2g-b$ gives
$g'=g+1$.  A $1$-handle also has two attaching disks and one lateral annulus,
which proves the last assertion.
\end{proof}

\begin{remark}[Extent of the surgery statement]
\label{rem:tunnel-scope}
A blind indentation replaces one boundary disk by another and adds no handle.
If the two mouths lie on distinct boundary components, the operation joins
those components: their genera add, rather than increasing one genus by one.
The proposition determines the boundary surface even for a knotted tunnel; it
does not classify the remaining $3$-manifold.  Removing a tunnel and attaching
material outside the body are also different operations on an embedded body.
The next product model fixes an unknotted tunnel and proves its body type
directly.
\end{remark}

\begin{proposition}[A pierced cap becomes an annulus]
\label{prop:pierced-cap-annulus}
Let $I=[0,1]$, let $D^2$ be an oriented closed round disk, and let $d^2$ be a
smaller concentric closed disk in its interior.  Set
\[
 H_0=I\times D^2,\qquad
 H_1=I\times\bigl(D^2\setminus\operatorname{int}d^2\bigr).
\]
Then $H_0$ is a $3$-ball and $H_1$ is a solid torus.  For $0<t_0<1$, put
\[
 D=\{t_0\}\times D^2,\quad C=\partial D,\quad
 d=\{t_0\}\times d^2,\quad c=\partial d,\quad
 A=D\setminus\operatorname{int}d.
\]
Give $c$ the boundary orientation of $d$, so that the inner boundary of $A$
has orientation $-c$.  As oriented chains,
\begin{equation}\label{eq:pierced-cap-chain}
 \partial D=C,\qquad \partial A=C-c.
\end{equation}
The same surviving curve $C$ satisfies
\[
 [C]=0\ \text{in }H_1(H_0;\mathbb Z),\qquad
 [C]=[c]\ne0\ \text{in }H_1(H_1;\mathbb Z)\cong\mathbb Z.
\]
In particular, $C$ bounds no disk, or integral $2$-chain, in the drilled body.
\end{proposition}

\begin{proof}
The annulus $D^2\setminus\operatorname{int}d^2$ is $S^1\times I$, hence
$H_1\cong S^1\times I^2\cong S^1\times D^2$ after rounding corners.
Subtracting the chain $d$ from $D$ gives
\eqref{eq:pierced-cap-chain}.  Retraction of the annulus onto its circular
factor, together with contraction of the first interval, is a deformation
retraction of $H_1$ onto $S^1$.  Both $C$ and $c$ map with degree one to this
circle, proving that their common class generates $H_1(H_1;\mathbb Z)$.
In $H_0$ the original disk $D$ fills $C$.
\end{proof}

\begin{remark}[Which circle still has a cap?]
\label{rem:drilled-body-meridian}
There are still disk-bounding curves on the new boundary.  If $r$ is a radial
arc joining the two boundaries of the annulus, then $I\times r$ is a properly
embedded disk in $H_1$.  Its boundary is a meridian of the remaining solid
torus, whereas $C$ represents its surviving longitudinal class.  In particular,
the inclusion $\partial H_1\hookrightarrow H_1$ has
\[
 H_1(\partial H_1;\mathbb Z)\cong\mathbb Z^2
 \longrightarrow H_1(H_1;\mathbb Z)\cong\mathbb Z
\]
with kernel generated by that meridian.  Calling a curve a ``hole boundary''
does not decide whether it has a filling disk; the body and the curve both
matter.  Nor, in general, is vanishing of a homology class by itself a
certificate of an embedded disk filling.
\end{remark}

\subsection{Double pants gluing and annular interface data}
\label{subsec:double-pants-annular-interface}

\begin{proposition}[A handle without directed feedback]
\label{prop:double-pants-dag-handle}
In the declared minimal thickening, join a $1\to2$ gate and a $2\to1$ gate
by two distinct output--input wires.  The side surface is $S_{1,2}$, obtained
by gluing two pants along two pairs of boundary circles.  The directed graph
is acyclic, while its underlying multigraph has first Betti number one.
\end{proposition}

\begin{proof}
Each pants has Euler characteristic $-1$; gluing along circles preserves their
sum $-2$, and four of the six boundary circles are consumed.  Thus $b=2$ and
$2-2g-2=-2$, giving $g=1$.  Both wires point from the first gate to the second,
so there is no directed cycle.  The underlying connected multigraph has two
vertices and two edges, giving $\beta_1=2-2+1=1$, in agreement with
Proposition~\ref{prop:local-pants-handlebody-genus}.  The first gluing connects
the two components; the second creates the handle.
\end{proof}

\begin{example}[The algebraic decoration is additional data]
\label{ex:split-merge-handle}
Declare a copy gate $\Delta(x)=(x,x)$, unary branch maps $f,g$, and a binary
merge $m$.  The corresponding computation is
\[
 x\longmapsto(f(x),g(x))\longmapsto m(f(x),g(x)),
\]
on its declared partial domain.  Unary gates only subdivide the two wires and
do not change the genus.  For $f=g=\id$ and $m(u,v)=u+v$ the output is $2x$,
so handle insertion does not by itself preserve the previous numerical
function.  The two parallel wires also show why the event reachability poset
alone is insufficient: one wire and two wires between the same two events
give the same order but different $\beta_1$.  Port incidence and multiplicity
must be retained alongside $P$ and $\mathcal O(P)$.
\end{example}

\begin{proposition}[Annular Stokes balance]
\label{prop:annular-stokes-balance}
In Proposition~\ref{prop:pierced-cap-annulus}, suppose a smooth one-form
$\widetilde\beta$ is declared on a neighborhood of $A$ in the drilled body.
Using the orientations in \eqref{eq:pierced-cap-chain},
\begin{equation}\label{eq:annular-stokes-balance}
 \oint_C\widetilde\beta-\oint_c\widetilde\beta
       =\int_A d\widetilde\beta.
\end{equation}
If $\widetilde\beta=-\lambda a\,dv$ for smooth declared data $a,v$ and constant
$\lambda$, this is the annular version of the seam-circle formula.  It requires
no extension across the removed disk.
\end{proposition}

\begin{proof}
Apply Stokes to the compact oriented annulus $A$ with boundary $C-c$.
For the specified potential, $d\widetilde\beta=-\lambda\,da\wedge dv$.
\end{proof}

\begin{remark}[Periods and the open arithmetic interface]
\label{rem:annular-interface-open}
If $d\widetilde\beta=0$, the two periods in
\eqref{eq:annular-stokes-balance} agree but need not vanish.  For example, on
the punctured disk with coordinates $(\xi,\eta)$,
\[
 \omega=\frac{-\eta\,d\xi+\xi\,d\eta}{\xi^2+\eta^2}
\]
is smooth and closed, and both positively oriented circular periods are
$2\pi$.  This is an example for a general one-form; it does not assert a
global representation of $\omega$ by arithmetic data $-\lambda a\,dv$.

Topological non-fillability neither forces a nonzero period nor assigns a
spectral value to a circle.  A proposed arithmetic annular interface must
specify both boundary state spaces and their compatibility relation, retaining
source, occurrence and production residuals.  Neither the topological surgery
nor \eqref{eq:annular-stokes-balance} supplies that relation, a regular AES
realization, or a native \texttt{Seal} rule.  These remain open under OQ-096;
the product-model calibration is recorded in
Appendix~\ref{app:pants-cellulation-verification}.
\end{remark}

\subsection{Interface state, substitution, and the spectrum}

\begin{proposition}[Polynomial interface spectrum]
\label{prop:polynomial-interface-spectrum}
In the rational-polynomial fragment, with $Q_+:z=x+y$ and $P_\times:w=zt$ glued
along the interface disk $D_z$, the substitution is the algebra homomorphism
\[
  s^*:\mathbb Q[z]\to\mathbb Q[x,y],\qquad z\mapsto x+y ,
\]
and evaluation characters satisfy $\chi_{r,s}\circ s^*=\chi_{r+s}$.  The
``spectrum value'' on the interface is the rational point $\chi_c(z)=c$ of
$\operatorname{Spec}\mathbb Q[z]$, and the relation gluing is the tensor product
(geometric fiber product)
\[
  \mathbb Q[x,y,z]/(z-x-y)\otimes_{\mathbb Q[z]}\mathbb Q[z,t,w]/(w-zt)
  \cong\mathbb Q[x,y,t,w]/(w-(x+y)t).
\]
The spectrum $\operatorname{Spec}A_z$ is a state-parameter space, not the
geometric seam circle $C_z$ and not the Stone order spectrum of
Theorem~\ref{thm:birkhoff-downset-spinal}.
\end{proposition}

\begin{proof}
Direct substitution and elimination of $z$; the choice of this model as the AEG
interface-spectrum semantics is a structural proposal, the displayed algebra
itself being a proved identity in the polynomial fragment.
\end{proof}

\begin{definition}[Interface record with residual]
\label{def:interface-record-with-residual}
An interface disk record is the five-tuple
\[
  \mathcal I_e=(D_e,C_e;\sigma_e,r_e,\beta_e),
\]
where $\sigma_e$ is the interface state under the declared observation layer,
$r_e$ is a recoverable production-process residual (together with any tangential
geometric residual), and $\beta_e$ is the port binding credential.  Numerical
equality of $\sigma_e$ does not identify two production histories ($2+3$ and
$1+4$ both carry $\sigma_e=5$).
\end{definition}

\subsection{The capping obstruction and certified sealing}

\begin{theorem}[Constant-boundary capping obstruction]
\label{thm:constant-boundary-capping-obstruction}
Let $D$ be a compact smooth disk, $g$ a smooth Riemannian metric, $a\in C^\infty(D)$,
and $\mu\ne0$, $\lambda\in\mathbb R$.  There is no data $(g,a)$ satisfying
$|da|_g^2=\mu^2+\lambda^2 a^2$ on $\operatorname{int}D$ and $a|_{\partial D}=c$
(constant).
\end{theorem}

\begin{proof}
Since $\mu^2+\lambda^2a^2\ge\mu^2>0$, one has $da\ne0$ throughout the interior.
If $a$ is nonconstant, compactness gives an interior extremum of $a$, where
$da=0$, a contradiction; if $a$ is constant, $da=0$ everywhere, also a
contradiction.
\end{proof}

\begin{corollary}[Capping dichotomy]
\label{cor:capping-dichotomy}
\begin{enumerate}
  \item An interface disk with numeric label $\sigma_e=c$ must not be realized
  as a compact regular AES disk with constant boundary assignment
  $a|_{C_e}=c$.
  \item To cap a level circle geometrically one must admit an interior critical
  point or a non-$C^1$/metric singularity (the singular category of
  Section~\ref{sec:zero-geometry}).
  \item To glue across a seam $C_e$ in the regular category, both sides must
  satisfy the two-sided smooth collar contract of Paper~0
  ($\texttt{def:p0-collar-contract}$).  Nonconstant $a|_{C_e}$ is necessary
  if a compact regular AES disk is also to fill that seam.  The disk
  obstruction alone imposes no such conclusion on gluing annular collars
  without a disk filling.
\end{enumerate}
\end{corollary}

\begin{definition}[Certified cap tuple]
\label{def:certified-cap-tuple}
A candidate \emph{certified cap} is a six-tuple
\[
  \mathcal C=(D,C,\sigma,r,w;\mathcal Q),
\]
where $D,C$ are the declared fill and its edge, $\sigma$ a state, $r$ a
residual, $w$ a verification certificate against an explicit closure obligation,
and $\mathcal Q$ the declared observation scope.  Four closure judgments are kept
distinct: (i) computational port substitution binding, (ii) topological
disk-filling $\partial D=C$ or conformal cusp filling, (iii) regular AES collar
descent or essential-singularity removal, and (iv) the Adva
\texttt{Seal \{ body \}} witness closure (additive residual zero, multiplicative
residual unit, with the non-vanishing domain obligations retained).  A
substitution success does not imply \texttt{Seal}, and $A=0,M=1$ does not by
itself produce an AES cap.
\end{definition}
